
\documentclass{article}
\usepackage{colortbl}
\usepackage{makecell}
\usepackage{multirow}
\usepackage{supertabular}

\begin{document}

\newcounter{utterance}

\centering \large Interaction Transcript for game `codenames', experiment `association\_difficult', episode 4 with gpt{-}5.4{-}azure.
\vspace{24pt}

{ \footnotesize  \setcounter{utterance}{1}
\setlength{\tabcolsep}{0pt}
\begin{supertabular}{c@{$\;$}|p{.15\linewidth}@{}p{.15\linewidth}p{.15\linewidth}p{.15\linewidth}p{.15\linewidth}p{.15\linewidth}}
    \# & $\;$A & \multicolumn{4}{c}{Game Master} & $\;\:$B\\
    \hline

    \theutterance \stepcounter{utterance}  
    & & & \multicolumn{2}{p{0.3\linewidth}}{
        \cellcolor[rgb]{0.9,0.9,0.9}{
            \makecell[{{p{\linewidth}}}]{
                \texttt{\tiny{[GM$|$GM]}}
                {'hidden': {'team': ['orhidee', 'jaanalind', 'nartsiss', 'ranne', 'peet', 'takso', 'rumm', 'jacks', 'jänes'], 'innocent': ['maastur', 'scotch', 'salat', 'veoauto', 'põder', 'pilkenär', 'reis'], 'opponent': ['kajakas', 'jalgratas', 'võõrasema', 'ale', 'sibul', 'traktor', 'tulp', 'maitseained'], 'assassin': ['varblane']}, 'revealed': {'team': {'team': [], 'innocent': [], 'opponent': [], 'assassin': []}, 'opponent': {'team': [], 'innocent': [], 'opponent': [], 'assassin': []}}}
            }
        }
    }
    & & \\ \\

    \theutterance \stepcounter{utterance}  
    & & \multicolumn{4}{p{0.6\linewidth}}{
        \cellcolor[rgb]{0.9,0.9,0.9}{
            \makecell[{{p{\linewidth}}}]{
                \texttt{\tiny{[P1$\langle$GM]}}
                \texttt{Esitage üks ühe sõnaga vihje, mis on seotud mõne teie meeskonna sõnaga allpool esitatud nimekirjast. Võite valida, kui paljude ja milliste sõnadega teie vihje seotud on.} \\
\texttt{Eesmärk on, et teie partner peab teie antud vihjete põhjal ära arvama kõik teie meeskonna sõnad, enne kui teine meeskond on kõik oma sõnad avastanud. Kui teie partner arvab ära "vastase" või "segaja" sõna, lõpeb nende arvamisvoor ja vastasmeeskond saab teha käigu. Kui teie partner arvab ära "tapa" sõna, lõpeb mäng ja teie meeskond kaotab kohe.} \\
\texttt{Mida rohkemate sõnadega teie vihje seostub, seda kiiremini võite mängu võita, kuid seda raskem võib olla teie partneril seoseid ära arvata, seega valige oma vihjesõna targalt. Vihjesõna peab olema semantiliselt seotud sihtsõnadega, see ei tohi olla üks nimekirjades olevatest sõnadest ega sisaldada nende osasid.} \\
\\ 
\texttt{Andke alati oma ühe sõnaga vihje ja oma komadega eraldatud sihtsõnade loetelu järgmises vormis ning tehke oma vastused võimalikult lühikeseks, ärge kunagi lisage muud teksti kui selles vormis nõutud:} \\
\\ 
\texttt{VIHJE: <SÕNA>} \\
\texttt{SIHID: <SIHID>} \\
\\ 
\texttt{Teie meeskonna sõnad on: orhidee, jaanalind, nartsiss, ranne, peet, takso, rumm, jacks, jänes.} \\
\texttt{Teie vastase meeskonna sõnad on: kajakas, jalgratas, võõrasema, ale, sibul, traktor, tulp, maitseained.} \\
\texttt{Segajate sõnad on: maastur, scotch, salat, veoauto, põder, pilkenär, reis.} \\
\texttt{Tapa sõnad on: varblane.} \\
            }
        }
    }
    & \\ \\

    \theutterance \stepcounter{utterance}  
    & \multicolumn{4}{p{0.6\linewidth}}{
        \cellcolor[rgb]{0.9,0.9,0.9}{
            \makecell[{{p{\linewidth}}}]{
                \texttt{\tiny{[P1$\rangle$GM]}}
                \texttt{VIHJE: edevus} \\
\texttt{SIHID: nartsiss} \\
            }
        }
    }
    & & \\ \\

    \theutterance \stepcounter{utterance}  
    & & & \multicolumn{2}{p{0.3\linewidth}}{
        \cellcolor[rgb]{0.9,0.9,0.9}{
            \makecell[{{p{\linewidth}}}]{
                \texttt{\tiny{[GM$|$GM]}}
                \texttt{edevus} \\
            }
        }
    }
    & & \\ \\

    \theutterance \stepcounter{utterance}  
    & & & \multicolumn{2}{p{0.3\linewidth}}{
        \cellcolor[rgb]{0.9,0.9,0.9}{
            \makecell[{{p{\linewidth}}}]{
                \texttt{\tiny{[GM$|$GM]}}
                ['nartsiss']
            }
        }
    }
    & & \\ \\

    \theutterance \stepcounter{utterance}  
    & & & \multicolumn{2}{p{0.3\linewidth}}{
        \cellcolor[rgb]{0.9,0.9,0.9}{
            \makecell[{{p{\linewidth}}}]{
                \texttt{\tiny{[GM$|$GM]}}
                {'word': 'nartsiss', 'assignment': 'team'}
            }
        }
    }
    & & \\ \\

    \theutterance \stepcounter{utterance}  
    & & \multicolumn{4}{p{0.6\linewidth}}{
        \cellcolor[rgb]{0.9,0.9,0.9}{
            \makecell[{{p{\linewidth}}}]{
                \texttt{\tiny{[GM$\rangle$P2]}}
                \texttt{Esitage komadega eraldatud loetelu kuni 1 sõnast järgmisest loetelust, mis on kõige paremini seotud või kõige tihedamalt seostuvad sõnaga 'edevus'. Alustage oma arvamuste loetelu alati sõnadega ARVING:  ja ärge lisage oma vastusesse muud teksti.} \\
\\ 
\texttt{põder, scotch, nartsiss, jaanalind, salat, maitseained, peet, pilkenär, traktor, maastur, veoauto, jalgratas, takso, rumm, tulp, ale, sibul, orhidee, varblane, ranne, jacks, jänes, võõrasema, kajakas, reis} \\
            }
        }
    }
    & \\ \\

    \theutterance \stepcounter{utterance}  
    & & & \multicolumn{4}{p{0.6\linewidth}}{
        \cellcolor[rgb]{0.9,0.9,0.9}{
            \makecell[{{p{\linewidth}}}]{
                \texttt{\tiny{[GM$\langle$P2]}}
                \texttt{ARVING: nartsiss} \\
            }
        }
    }
     \\ \\

    \theutterance \stepcounter{utterance}  
    & & & \multicolumn{2}{p{0.3\linewidth}}{
        \cellcolor[rgb]{0.9,0.9,0.9}{
            \makecell[{{p{\linewidth}}}]{
                \texttt{\tiny{[GM$|$GM]}}
                ['nartsiss']
            }
        }
    }
    & & \\ \\

    \theutterance \stepcounter{utterance}  
    & & & \multicolumn{2}{p{0.3\linewidth}}{
        \cellcolor[rgb]{0.9,0.9,0.9}{
            \makecell[{{p{\linewidth}}}]{
                \texttt{\tiny{[GM$|$GM]}}
                {'word': 'nartsiss', 'assignment': 'team'}
            }
        }
    }
    & & \\ \\

    \theutterance \stepcounter{utterance}  
    & & & \multicolumn{2}{p{0.3\linewidth}}{
        \cellcolor[rgb]{0.9,0.9,0.9}{
            \makecell[{{p{\linewidth}}}]{
                \texttt{\tiny{[GM$|$GM]}}
                {'word': 'nartsiss', 'assignment': 'team'}
            }
        }
    }
    & & \\ \\

    \theutterance \stepcounter{utterance}  
    & & & \multicolumn{2}{p{0.3\linewidth}}{
        \cellcolor[rgb]{0.9,0.9,0.9}{
            \makecell[{{p{\linewidth}}}]{
                \texttt{\tiny{[GM$|$GM]}}
                {'word': 'tulp', 'assignment': 'opponent'}
            }
        }
    }
    & & \\ \\

    \theutterance \stepcounter{utterance}  
    & & & \multicolumn{2}{p{0.3\linewidth}}{
        \cellcolor[rgb]{0.9,0.9,0.9}{
            \makecell[{{p{\linewidth}}}]{
                \texttt{\tiny{[GM$|$GM]}}
                {'hidden': {'team': ['orhidee', 'jaanalind', 'ranne', 'peet', 'takso', 'rumm', 'jacks', 'jänes'], 'innocent': ['maastur', 'scotch', 'salat', 'veoauto', 'põder', 'pilkenär', 'reis'], 'opponent': ['kajakas', 'jalgratas', 'võõrasema', 'ale', 'sibul', 'traktor', 'maitseained'], 'assassin': ['varblane']}, 'revealed': {'team': {'team': ['nartsiss'], 'innocent': [], 'opponent': [], 'assassin': []}, 'opponent': {'team': [], 'innocent': [], 'opponent': ['tulp'], 'assassin': []}}}
            }
        }
    }
    & & \\ \\

    \theutterance \stepcounter{utterance}  
    & & \multicolumn{4}{p{0.6\linewidth}}{
        \cellcolor[rgb]{0.9,0.9,0.9}{
            \makecell[{{p{\linewidth}}}]{
                \texttt{\tiny{[P1$\langle$GM]}}
                \texttt{Sõna nartsiss arvati õigesti. Teie meeskonnakaaslase käik lõppes.} \\
\texttt{Nüüd andke veel üks vihje, mis on seotud mõne teie ülejäänud meeskonnasõnaga, ja nimekiri seotud sihtmärgiks olevatest sõnadest. Ärge unustage alustada oma vihjet sõnaga 'VIHJE: ', pange uus rida ja alustage oma komadega eraldatud sihtmärgisõnade loetelu sõnaga 'SIHID: '. Märkus: mõned sõnad on võrreldes eelmiste taotlustega nimekirjadest eemaldatud.} \\
\\ 
\texttt{Teie ülejäänud meeskonnasõnad on: orhidee, jaanalind, ranne, peet, takso, rumm, jacks, jänes.} \\
\texttt{Teie vastase ülejäänud sõnad on: kajakas, jalgratas, võõrasema, ale, sibul, traktor, tulp, maitseained.} \\
\texttt{Ülejäänud segavad sõnad on: maastur, scotch, salat, veoauto, põder, pilkenär, reis.} \\
\texttt{Ülejäänud tapmissõnad on: varblane.} \\
            }
        }
    }
    & \\ \\

    \theutterance \stepcounter{utterance}  
    & \multicolumn{4}{p{0.6\linewidth}}{
        \cellcolor[rgb]{0.9,0.9,0.9}{
            \makecell[{{p{\linewidth}}}]{
                \texttt{\tiny{[P1$\rangle$GM]}}
                \texttt{VIHJE: jänku} \\
\texttt{SIHID: jänes} \\
            }
        }
    }
    & & \\ \\

    \theutterance \stepcounter{utterance}  
    & & & \multicolumn{2}{p{0.3\linewidth}}{
        \cellcolor[rgb]{0.9,0.9,0.9}{
            \makecell[{{p{\linewidth}}}]{
                \texttt{\tiny{[GM$|$GM]}}
                \texttt{jänku} \\
            }
        }
    }
    & & \\ \\

    \theutterance \stepcounter{utterance}  
    & & & \multicolumn{2}{p{0.3\linewidth}}{
        \cellcolor[rgb]{0.9,0.9,0.9}{
            \makecell[{{p{\linewidth}}}]{
                \texttt{\tiny{[GM$|$GM]}}
                ['jänes']
            }
        }
    }
    & & \\ \\

    \theutterance \stepcounter{utterance}  
    & & & \multicolumn{2}{p{0.3\linewidth}}{
        \cellcolor[rgb]{0.9,0.9,0.9}{
            \makecell[{{p{\linewidth}}}]{
                \texttt{\tiny{[GM$|$GM]}}
                {'word': 'jänes', 'assignment': 'team'}
            }
        }
    }
    & & \\ \\

    \theutterance \stepcounter{utterance}  
    & & \multicolumn{4}{p{0.6\linewidth}}{
        \cellcolor[rgb]{0.9,0.9,0.9}{
            \makecell[{{p{\linewidth}}}]{
                \texttt{\tiny{[GM$\rangle$P2]}}
                \texttt{Sõna nartsiss arvati õigesti. Teie käik lõppes.} \\
\texttt{Nüüd esitage veel üks komadega eraldatud loetelu vähemalt 1 ja kuni 1 sõnast järgmisest sõnade loetelust, mis on kõige paremini seotud või kõige tihedamalt seotud sõnaga 'jänku'. Ärge unustage alustada oma vastust sõnadega ARVING: . Märkus: mõned sõnad on võrreldes eelmiste päringutega nimekirjast eemaldatud.} \\
\\ 
\texttt{põder, scotch, jaanalind, salat, maitseained, peet, pilkenär, traktor, maastur, veoauto, jalgratas, takso, rumm, ale, sibul, orhidee, varblane, ranne, jacks, jänes, võõrasema, kajakas, reis} \\
            }
        }
    }
    & \\ \\

    \theutterance \stepcounter{utterance}  
    & & & \multicolumn{4}{p{0.6\linewidth}}{
        \cellcolor[rgb]{0.9,0.9,0.9}{
            \makecell[{{p{\linewidth}}}]{
                \texttt{\tiny{[GM$\langle$P2]}}
                \texttt{ARVING: jänes} \\
            }
        }
    }
     \\ \\

    \theutterance \stepcounter{utterance}  
    & & & \multicolumn{2}{p{0.3\linewidth}}{
        \cellcolor[rgb]{0.9,0.9,0.9}{
            \makecell[{{p{\linewidth}}}]{
                \texttt{\tiny{[GM$|$GM]}}
                ['jänes']
            }
        }
    }
    & & \\ \\

    \theutterance \stepcounter{utterance}  
    & & & \multicolumn{2}{p{0.3\linewidth}}{
        \cellcolor[rgb]{0.9,0.9,0.9}{
            \makecell[{{p{\linewidth}}}]{
                \texttt{\tiny{[GM$|$GM]}}
                {'word': 'jänes', 'assignment': 'team'}
            }
        }
    }
    & & \\ \\

    \theutterance \stepcounter{utterance}  
    & & & \multicolumn{2}{p{0.3\linewidth}}{
        \cellcolor[rgb]{0.9,0.9,0.9}{
            \makecell[{{p{\linewidth}}}]{
                \texttt{\tiny{[GM$|$GM]}}
                {'word': 'jänes', 'assignment': 'team'}
            }
        }
    }
    & & \\ \\

    \theutterance \stepcounter{utterance}  
    & & & \multicolumn{2}{p{0.3\linewidth}}{
        \cellcolor[rgb]{0.9,0.9,0.9}{
            \makecell[{{p{\linewidth}}}]{
                \texttt{\tiny{[GM$|$GM]}}
                {'word': 'traktor', 'assignment': 'opponent'}
            }
        }
    }
    & & \\ \\

    \theutterance \stepcounter{utterance}  
    & & & \multicolumn{2}{p{0.3\linewidth}}{
        \cellcolor[rgb]{0.9,0.9,0.9}{
            \makecell[{{p{\linewidth}}}]{
                \texttt{\tiny{[GM$|$GM]}}
                {'hidden': {'team': ['orhidee', 'jaanalind', 'ranne', 'peet', 'takso', 'rumm', 'jacks'], 'innocent': ['maastur', 'scotch', 'salat', 'veoauto', 'põder', 'pilkenär', 'reis'], 'opponent': ['kajakas', 'jalgratas', 'võõrasema', 'ale', 'sibul', 'maitseained'], 'assassin': ['varblane']}, 'revealed': {'team': {'team': ['nartsiss', 'jänes'], 'innocent': [], 'opponent': [], 'assassin': []}, 'opponent': {'team': [], 'innocent': [], 'opponent': ['tulp', 'traktor'], 'assassin': []}}}
            }
        }
    }
    & & \\ \\

    \theutterance \stepcounter{utterance}  
    & & \multicolumn{4}{p{0.6\linewidth}}{
        \cellcolor[rgb]{0.9,0.9,0.9}{
            \makecell[{{p{\linewidth}}}]{
                \texttt{\tiny{[P1$\langle$GM]}}
                \texttt{Sõna jänes arvati õigesti. Teie meeskonnakaaslase käik lõppes.} \\
\texttt{Nüüd andke veel üks vihje, mis on seotud mõne teie ülejäänud meeskonnasõnaga, ja nimekiri seotud sihtmärgiks olevatest sõnadest. Ärge unustage alustada oma vihjet sõnaga 'VIHJE: ', pange uus rida ja alustage oma komadega eraldatud sihtmärgisõnade loetelu sõnaga 'SIHID: '. Märkus: mõned sõnad on võrreldes eelmiste taotlustega nimekirjadest eemaldatud.} \\
\\ 
\texttt{Teie ülejäänud meeskonnasõnad on: orhidee, jaanalind, ranne, peet, takso, rumm, jacks.} \\
\texttt{Teie vastase ülejäänud sõnad on: kajakas, jalgratas, võõrasema, ale, sibul, traktor, maitseained.} \\
\texttt{Ülejäänud segavad sõnad on: maastur, scotch, salat, veoauto, põder, pilkenär, reis.} \\
\texttt{Ülejäänud tapmissõnad on: varblane.} \\
            }
        }
    }
    & \\ \\

    \theutterance \stepcounter{utterance}  
    & \multicolumn{4}{p{0.6\linewidth}}{
        \cellcolor[rgb]{0.9,0.9,0.9}{
            \makecell[{{p{\linewidth}}}]{
                \texttt{\tiny{[P1$\rangle$GM]}}
                \texttt{VIHJE: piraat} \\
\texttt{SIHID: rumm} \\
            }
        }
    }
    & & \\ \\

    \theutterance \stepcounter{utterance}  
    & & & \multicolumn{2}{p{0.3\linewidth}}{
        \cellcolor[rgb]{0.9,0.9,0.9}{
            \makecell[{{p{\linewidth}}}]{
                \texttt{\tiny{[GM$|$GM]}}
                \texttt{piraat} \\
            }
        }
    }
    & & \\ \\

    \theutterance \stepcounter{utterance}  
    & & & \multicolumn{2}{p{0.3\linewidth}}{
        \cellcolor[rgb]{0.9,0.9,0.9}{
            \makecell[{{p{\linewidth}}}]{
                \texttt{\tiny{[GM$|$GM]}}
                ['rumm']
            }
        }
    }
    & & \\ \\

    \theutterance \stepcounter{utterance}  
    & & & \multicolumn{2}{p{0.3\linewidth}}{
        \cellcolor[rgb]{0.9,0.9,0.9}{
            \makecell[{{p{\linewidth}}}]{
                \texttt{\tiny{[GM$|$GM]}}
                {'word': 'rumm', 'assignment': 'team'}
            }
        }
    }
    & & \\ \\

    \theutterance \stepcounter{utterance}  
    & & \multicolumn{4}{p{0.6\linewidth}}{
        \cellcolor[rgb]{0.9,0.9,0.9}{
            \makecell[{{p{\linewidth}}}]{
                \texttt{\tiny{[GM$\rangle$P2]}}
                \texttt{Sõna jänes arvati õigesti. Teie käik lõppes.} \\
\texttt{Nüüd esitage veel üks komadega eraldatud loetelu vähemalt 1 ja kuni 1 sõnast järgmisest sõnade loetelust, mis on kõige paremini seotud või kõige tihedamalt seotud sõnaga 'piraat'. Ärge unustage alustada oma vastust sõnadega ARVING: . Märkus: mõned sõnad on võrreldes eelmiste päringutega nimekirjast eemaldatud.} \\
\\ 
\texttt{põder, scotch, jaanalind, salat, maitseained, peet, pilkenär, maastur, veoauto, jalgratas, takso, rumm, ale, sibul, orhidee, varblane, ranne, jacks, võõrasema, kajakas, reis} \\
            }
        }
    }
    & \\ \\

    \theutterance \stepcounter{utterance}  
    & & & \multicolumn{4}{p{0.6\linewidth}}{
        \cellcolor[rgb]{0.9,0.9,0.9}{
            \makecell[{{p{\linewidth}}}]{
                \texttt{\tiny{[GM$\langle$P2]}}
                \texttt{ARVING: rumm} \\
            }
        }
    }
     \\ \\

    \theutterance \stepcounter{utterance}  
    & & & \multicolumn{2}{p{0.3\linewidth}}{
        \cellcolor[rgb]{0.9,0.9,0.9}{
            \makecell[{{p{\linewidth}}}]{
                \texttt{\tiny{[GM$|$GM]}}
                ['rumm']
            }
        }
    }
    & & \\ \\

    \theutterance \stepcounter{utterance}  
    & & & \multicolumn{2}{p{0.3\linewidth}}{
        \cellcolor[rgb]{0.9,0.9,0.9}{
            \makecell[{{p{\linewidth}}}]{
                \texttt{\tiny{[GM$|$GM]}}
                {'word': 'rumm', 'assignment': 'team'}
            }
        }
    }
    & & \\ \\

    \theutterance \stepcounter{utterance}  
    & & & \multicolumn{2}{p{0.3\linewidth}}{
        \cellcolor[rgb]{0.9,0.9,0.9}{
            \makecell[{{p{\linewidth}}}]{
                \texttt{\tiny{[GM$|$GM]}}
                {'word': 'rumm', 'assignment': 'team'}
            }
        }
    }
    & & \\ \\

    \theutterance \stepcounter{utterance}  
    & & & \multicolumn{2}{p{0.3\linewidth}}{
        \cellcolor[rgb]{0.9,0.9,0.9}{
            \makecell[{{p{\linewidth}}}]{
                \texttt{\tiny{[GM$|$GM]}}
                {'word': 'ale', 'assignment': 'opponent'}
            }
        }
    }
    & & \\ \\

    \theutterance \stepcounter{utterance}  
    & & & \multicolumn{2}{p{0.3\linewidth}}{
        \cellcolor[rgb]{0.9,0.9,0.9}{
            \makecell[{{p{\linewidth}}}]{
                \texttt{\tiny{[GM$|$GM]}}
                {'hidden': {'team': ['orhidee', 'jaanalind', 'ranne', 'peet', 'takso', 'jacks'], 'innocent': ['maastur', 'scotch', 'salat', 'veoauto', 'põder', 'pilkenär', 'reis'], 'opponent': ['kajakas', 'jalgratas', 'võõrasema', 'sibul', 'maitseained'], 'assassin': ['varblane']}, 'revealed': {'team': {'team': ['nartsiss', 'jänes', 'rumm'], 'innocent': [], 'opponent': [], 'assassin': []}, 'opponent': {'team': [], 'innocent': [], 'opponent': ['tulp', 'traktor', 'ale'], 'assassin': []}}}
            }
        }
    }
    & & \\ \\

    \theutterance \stepcounter{utterance}  
    & & \multicolumn{4}{p{0.6\linewidth}}{
        \cellcolor[rgb]{0.9,0.9,0.9}{
            \makecell[{{p{\linewidth}}}]{
                \texttt{\tiny{[P1$\langle$GM]}}
                \texttt{Sõna rumm arvati õigesti. Teie meeskonnakaaslase käik lõppes.} \\
\texttt{Nüüd andke veel üks vihje, mis on seotud mõne teie ülejäänud meeskonnasõnaga, ja nimekiri seotud sihtmärgiks olevatest sõnadest. Ärge unustage alustada oma vihjet sõnaga 'VIHJE: ', pange uus rida ja alustage oma komadega eraldatud sihtmärgisõnade loetelu sõnaga 'SIHID: '. Märkus: mõned sõnad on võrreldes eelmiste taotlustega nimekirjadest eemaldatud.} \\
\\ 
\texttt{Teie ülejäänud meeskonnasõnad on: orhidee, jaanalind, ranne, peet, takso, jacks.} \\
\texttt{Teie vastase ülejäänud sõnad on: kajakas, jalgratas, võõrasema, ale, sibul, maitseained.} \\
\texttt{Ülejäänud segavad sõnad on: maastur, scotch, salat, veoauto, põder, pilkenär, reis.} \\
\texttt{Ülejäänud tapmissõnad on: varblane.} \\
            }
        }
    }
    & \\ \\

    \theutterance \stepcounter{utterance}  
    & \multicolumn{4}{p{0.6\linewidth}}{
        \cellcolor[rgb]{0.9,0.9,0.9}{
            \makecell[{{p{\linewidth}}}]{
                \texttt{\tiny{[P1$\rangle$GM]}}
                \texttt{VIHJE: käevõru} \\
\texttt{SIHID: ranne} \\
            }
        }
    }
    & & \\ \\

    \theutterance \stepcounter{utterance}  
    & & & \multicolumn{2}{p{0.3\linewidth}}{
        \cellcolor[rgb]{0.9,0.9,0.9}{
            \makecell[{{p{\linewidth}}}]{
                \texttt{\tiny{[GM$|$GM]}}
                \texttt{käevõru} \\
            }
        }
    }
    & & \\ \\

    \theutterance \stepcounter{utterance}  
    & & & \multicolumn{2}{p{0.3\linewidth}}{
        \cellcolor[rgb]{0.9,0.9,0.9}{
            \makecell[{{p{\linewidth}}}]{
                \texttt{\tiny{[GM$|$GM]}}
                ['ranne']
            }
        }
    }
    & & \\ \\

    \theutterance \stepcounter{utterance}  
    & & & \multicolumn{2}{p{0.3\linewidth}}{
        \cellcolor[rgb]{0.9,0.9,0.9}{
            \makecell[{{p{\linewidth}}}]{
                \texttt{\tiny{[GM$|$GM]}}
                {'word': 'ranne', 'assignment': 'team'}
            }
        }
    }
    & & \\ \\

    \theutterance \stepcounter{utterance}  
    & & \multicolumn{4}{p{0.6\linewidth}}{
        \cellcolor[rgb]{0.9,0.9,0.9}{
            \makecell[{{p{\linewidth}}}]{
                \texttt{\tiny{[GM$\rangle$P2]}}
                \texttt{Sõna rumm arvati õigesti. Teie käik lõppes.} \\
\texttt{Nüüd esitage veel üks komadega eraldatud loetelu vähemalt 1 ja kuni 1 sõnast järgmisest sõnade loetelust, mis on kõige paremini seotud või kõige tihedamalt seotud sõnaga 'käevõru'. Ärge unustage alustada oma vastust sõnadega ARVING: . Märkus: mõned sõnad on võrreldes eelmiste päringutega nimekirjast eemaldatud.} \\
\\ 
\texttt{põder, scotch, jaanalind, salat, maitseained, peet, pilkenär, maastur, veoauto, jalgratas, takso, sibul, orhidee, varblane, ranne, jacks, võõrasema, kajakas, reis} \\
            }
        }
    }
    & \\ \\

    \theutterance \stepcounter{utterance}  
    & & & \multicolumn{4}{p{0.6\linewidth}}{
        \cellcolor[rgb]{0.9,0.9,0.9}{
            \makecell[{{p{\linewidth}}}]{
                \texttt{\tiny{[GM$\langle$P2]}}
                \texttt{ARVING: ranne} \\
            }
        }
    }
     \\ \\

    \theutterance \stepcounter{utterance}  
    & & & \multicolumn{2}{p{0.3\linewidth}}{
        \cellcolor[rgb]{0.9,0.9,0.9}{
            \makecell[{{p{\linewidth}}}]{
                \texttt{\tiny{[GM$|$GM]}}
                ['ranne']
            }
        }
    }
    & & \\ \\

    \theutterance \stepcounter{utterance}  
    & & & \multicolumn{2}{p{0.3\linewidth}}{
        \cellcolor[rgb]{0.9,0.9,0.9}{
            \makecell[{{p{\linewidth}}}]{
                \texttt{\tiny{[GM$|$GM]}}
                {'word': 'ranne', 'assignment': 'team'}
            }
        }
    }
    & & \\ \\

    \theutterance \stepcounter{utterance}  
    & & & \multicolumn{2}{p{0.3\linewidth}}{
        \cellcolor[rgb]{0.9,0.9,0.9}{
            \makecell[{{p{\linewidth}}}]{
                \texttt{\tiny{[GM$|$GM]}}
                {'word': 'ranne', 'assignment': 'team'}
            }
        }
    }
    & & \\ \\

    \theutterance \stepcounter{utterance}  
    & & & \multicolumn{2}{p{0.3\linewidth}}{
        \cellcolor[rgb]{0.9,0.9,0.9}{
            \makecell[{{p{\linewidth}}}]{
                \texttt{\tiny{[GM$|$GM]}}
                {'word': 'jalgratas', 'assignment': 'opponent'}
            }
        }
    }
    & & \\ \\

    \theutterance \stepcounter{utterance}  
    & & & \multicolumn{2}{p{0.3\linewidth}}{
        \cellcolor[rgb]{0.9,0.9,0.9}{
            \makecell[{{p{\linewidth}}}]{
                \texttt{\tiny{[GM$|$GM]}}
                {'hidden': {'team': ['orhidee', 'jaanalind', 'peet', 'takso', 'jacks'], 'innocent': ['maastur', 'scotch', 'salat', 'veoauto', 'põder', 'pilkenär', 'reis'], 'opponent': ['kajakas', 'võõrasema', 'sibul', 'maitseained'], 'assassin': ['varblane']}, 'revealed': {'team': {'team': ['nartsiss', 'jänes', 'rumm', 'ranne'], 'innocent': [], 'opponent': [], 'assassin': []}, 'opponent': {'team': [], 'innocent': [], 'opponent': ['tulp', 'traktor', 'ale', 'jalgratas'], 'assassin': []}}}
            }
        }
    }
    & & \\ \\

    \theutterance \stepcounter{utterance}  
    & & \multicolumn{4}{p{0.6\linewidth}}{
        \cellcolor[rgb]{0.9,0.9,0.9}{
            \makecell[{{p{\linewidth}}}]{
                \texttt{\tiny{[P1$\langle$GM]}}
                \texttt{Sõna ranne arvati õigesti. Teie meeskonnakaaslase käik lõppes.} \\
\texttt{Nüüd andke veel üks vihje, mis on seotud mõne teie ülejäänud meeskonnasõnaga, ja nimekiri seotud sihtmärgiks olevatest sõnadest. Ärge unustage alustada oma vihjet sõnaga 'VIHJE: ', pange uus rida ja alustage oma komadega eraldatud sihtmärgisõnade loetelu sõnaga 'SIHID: '. Märkus: mõned sõnad on võrreldes eelmiste taotlustega nimekirjadest eemaldatud.} \\
\\ 
\texttt{Teie ülejäänud meeskonnasõnad on: orhidee, jaanalind, peet, takso, jacks.} \\
\texttt{Teie vastase ülejäänud sõnad on: kajakas, jalgratas, võõrasema, sibul, maitseained.} \\
\texttt{Ülejäänud segavad sõnad on: maastur, scotch, salat, veoauto, põder, pilkenär, reis.} \\
\texttt{Ülejäänud tapmissõnad on: varblane.} \\
            }
        }
    }
    & \\ \\

    \theutterance \stepcounter{utterance}  
    & \multicolumn{4}{p{0.6\linewidth}}{
        \cellcolor[rgb]{0.9,0.9,0.9}{
            \makecell[{{p{\linewidth}}}]{
                \texttt{\tiny{[P1$\rangle$GM]}}
                \texttt{VIHJE: borš} \\
\texttt{SIHID: peet} \\
            }
        }
    }
    & & \\ \\

    \theutterance \stepcounter{utterance}  
    & & & \multicolumn{2}{p{0.3\linewidth}}{
        \cellcolor[rgb]{0.9,0.9,0.9}{
            \makecell[{{p{\linewidth}}}]{
                \texttt{\tiny{[GM$|$GM]}}
                \texttt{borš} \\
            }
        }
    }
    & & \\ \\

    \theutterance \stepcounter{utterance}  
    & & & \multicolumn{2}{p{0.3\linewidth}}{
        \cellcolor[rgb]{0.9,0.9,0.9}{
            \makecell[{{p{\linewidth}}}]{
                \texttt{\tiny{[GM$|$GM]}}
                ['peet']
            }
        }
    }
    & & \\ \\

    \theutterance \stepcounter{utterance}  
    & & & \multicolumn{2}{p{0.3\linewidth}}{
        \cellcolor[rgb]{0.9,0.9,0.9}{
            \makecell[{{p{\linewidth}}}]{
                \texttt{\tiny{[GM$|$GM]}}
                {'word': 'peet', 'assignment': 'team'}
            }
        }
    }
    & & \\ \\

    \theutterance \stepcounter{utterance}  
    & & \multicolumn{4}{p{0.6\linewidth}}{
        \cellcolor[rgb]{0.9,0.9,0.9}{
            \makecell[{{p{\linewidth}}}]{
                \texttt{\tiny{[GM$\rangle$P2]}}
                \texttt{Sõna ranne arvati õigesti. Teie käik lõppes.} \\
\texttt{Nüüd esitage veel üks komadega eraldatud loetelu vähemalt 1 ja kuni 1 sõnast järgmisest sõnade loetelust, mis on kõige paremini seotud või kõige tihedamalt seotud sõnaga 'borš'. Ärge unustage alustada oma vastust sõnadega ARVING: . Märkus: mõned sõnad on võrreldes eelmiste päringutega nimekirjast eemaldatud.} \\
\\ 
\texttt{põder, scotch, jaanalind, salat, maitseained, peet, pilkenär, maastur, veoauto, takso, sibul, orhidee, varblane, jacks, võõrasema, kajakas, reis} \\
            }
        }
    }
    & \\ \\

    \theutterance \stepcounter{utterance}  
    & & & \multicolumn{4}{p{0.6\linewidth}}{
        \cellcolor[rgb]{0.9,0.9,0.9}{
            \makecell[{{p{\linewidth}}}]{
                \texttt{\tiny{[GM$\langle$P2]}}
                \texttt{ARVING: peet} \\
            }
        }
    }
     \\ \\

    \theutterance \stepcounter{utterance}  
    & & & \multicolumn{2}{p{0.3\linewidth}}{
        \cellcolor[rgb]{0.9,0.9,0.9}{
            \makecell[{{p{\linewidth}}}]{
                \texttt{\tiny{[GM$|$GM]}}
                ['peet']
            }
        }
    }
    & & \\ \\

    \theutterance \stepcounter{utterance}  
    & & & \multicolumn{2}{p{0.3\linewidth}}{
        \cellcolor[rgb]{0.9,0.9,0.9}{
            \makecell[{{p{\linewidth}}}]{
                \texttt{\tiny{[GM$|$GM]}}
                {'word': 'peet', 'assignment': 'team'}
            }
        }
    }
    & & \\ \\

    \theutterance \stepcounter{utterance}  
    & & & \multicolumn{2}{p{0.3\linewidth}}{
        \cellcolor[rgb]{0.9,0.9,0.9}{
            \makecell[{{p{\linewidth}}}]{
                \texttt{\tiny{[GM$|$GM]}}
                {'word': 'peet', 'assignment': 'team'}
            }
        }
    }
    & & \\ \\

    \theutterance \stepcounter{utterance}  
    & & & \multicolumn{2}{p{0.3\linewidth}}{
        \cellcolor[rgb]{0.9,0.9,0.9}{
            \makecell[{{p{\linewidth}}}]{
                \texttt{\tiny{[GM$|$GM]}}
                {'word': 'sibul', 'assignment': 'opponent'}
            }
        }
    }
    & & \\ \\

    \theutterance \stepcounter{utterance}  
    & & & \multicolumn{2}{p{0.3\linewidth}}{
        \cellcolor[rgb]{0.9,0.9,0.9}{
            \makecell[{{p{\linewidth}}}]{
                \texttt{\tiny{[GM$|$GM]}}
                {'hidden': {'team': ['orhidee', 'jaanalind', 'takso', 'jacks'], 'innocent': ['maastur', 'scotch', 'salat', 'veoauto', 'põder', 'pilkenär', 'reis'], 'opponent': ['kajakas', 'võõrasema', 'maitseained'], 'assassin': ['varblane']}, 'revealed': {'team': {'team': ['nartsiss', 'jänes', 'rumm', 'ranne', 'peet'], 'innocent': [], 'opponent': [], 'assassin': []}, 'opponent': {'team': [], 'innocent': [], 'opponent': ['tulp', 'traktor', 'ale', 'jalgratas', 'sibul'], 'assassin': []}}}
            }
        }
    }
    & & \\ \\

    \theutterance \stepcounter{utterance}  
    & & \multicolumn{4}{p{0.6\linewidth}}{
        \cellcolor[rgb]{0.9,0.9,0.9}{
            \makecell[{{p{\linewidth}}}]{
                \texttt{\tiny{[P1$\langle$GM]}}
                \texttt{Sõna peet arvati õigesti. Teie meeskonnakaaslase käik lõppes.} \\
\texttt{Nüüd andke veel üks vihje, mis on seotud mõne teie ülejäänud meeskonnasõnaga, ja nimekiri seotud sihtmärgiks olevatest sõnadest. Ärge unustage alustada oma vihjet sõnaga 'VIHJE: ', pange uus rida ja alustage oma komadega eraldatud sihtmärgisõnade loetelu sõnaga 'SIHID: '. Märkus: mõned sõnad on võrreldes eelmiste taotlustega nimekirjadest eemaldatud.} \\
\\ 
\texttt{Teie ülejäänud meeskonnasõnad on: orhidee, jaanalind, takso, jacks.} \\
\texttt{Teie vastase ülejäänud sõnad on: kajakas, võõrasema, sibul, maitseained.} \\
\texttt{Ülejäänud segavad sõnad on: maastur, scotch, salat, veoauto, põder, pilkenär, reis.} \\
\texttt{Ülejäänud tapmissõnad on: varblane.} \\
            }
        }
    }
    & \\ \\

    \theutterance \stepcounter{utterance}  
    & \multicolumn{4}{p{0.6\linewidth}}{
        \cellcolor[rgb]{0.9,0.9,0.9}{
            \makecell[{{p{\linewidth}}}]{
                \texttt{\tiny{[P1$\rangle$GM]}}
                \texttt{VIHJE: Uber} \\
\texttt{SIHID: takso} \\
            }
        }
    }
    & & \\ \\

    \theutterance \stepcounter{utterance}  
    & & & \multicolumn{2}{p{0.3\linewidth}}{
        \cellcolor[rgb]{0.9,0.9,0.9}{
            \makecell[{{p{\linewidth}}}]{
                \texttt{\tiny{[GM$|$GM]}}
                \texttt{uber} \\
            }
        }
    }
    & & \\ \\

    \theutterance \stepcounter{utterance}  
    & & & \multicolumn{2}{p{0.3\linewidth}}{
        \cellcolor[rgb]{0.9,0.9,0.9}{
            \makecell[{{p{\linewidth}}}]{
                \texttt{\tiny{[GM$|$GM]}}
                ['takso']
            }
        }
    }
    & & \\ \\

    \theutterance \stepcounter{utterance}  
    & & & \multicolumn{2}{p{0.3\linewidth}}{
        \cellcolor[rgb]{0.9,0.9,0.9}{
            \makecell[{{p{\linewidth}}}]{
                \texttt{\tiny{[GM$|$GM]}}
                {'word': 'takso', 'assignment': 'team'}
            }
        }
    }
    & & \\ \\

    \theutterance \stepcounter{utterance}  
    & & \multicolumn{4}{p{0.6\linewidth}}{
        \cellcolor[rgb]{0.9,0.9,0.9}{
            \makecell[{{p{\linewidth}}}]{
                \texttt{\tiny{[GM$\rangle$P2]}}
                \texttt{Sõna peet arvati õigesti. Teie käik lõppes.} \\
\texttt{Nüüd esitage veel üks komadega eraldatud loetelu vähemalt 1 ja kuni 1 sõnast järgmisest sõnade loetelust, mis on kõige paremini seotud või kõige tihedamalt seotud sõnaga 'uber'. Ärge unustage alustada oma vastust sõnadega ARVING: . Märkus: mõned sõnad on võrreldes eelmiste päringutega nimekirjast eemaldatud.} \\
\\ 
\texttt{põder, scotch, jaanalind, salat, maitseained, pilkenär, maastur, veoauto, takso, orhidee, varblane, jacks, võõrasema, kajakas, reis} \\
            }
        }
    }
    & \\ \\

    \theutterance \stepcounter{utterance}  
    & & & \multicolumn{4}{p{0.6\linewidth}}{
        \cellcolor[rgb]{0.9,0.9,0.9}{
            \makecell[{{p{\linewidth}}}]{
                \texttt{\tiny{[GM$\langle$P2]}}
                \texttt{ARVING: takso} \\
            }
        }
    }
     \\ \\

    \theutterance \stepcounter{utterance}  
    & & & \multicolumn{2}{p{0.3\linewidth}}{
        \cellcolor[rgb]{0.9,0.9,0.9}{
            \makecell[{{p{\linewidth}}}]{
                \texttt{\tiny{[GM$|$GM]}}
                ['takso']
            }
        }
    }
    & & \\ \\

    \theutterance \stepcounter{utterance}  
    & & & \multicolumn{2}{p{0.3\linewidth}}{
        \cellcolor[rgb]{0.9,0.9,0.9}{
            \makecell[{{p{\linewidth}}}]{
                \texttt{\tiny{[GM$|$GM]}}
                {'word': 'takso', 'assignment': 'team'}
            }
        }
    }
    & & \\ \\

    \theutterance \stepcounter{utterance}  
    & & & \multicolumn{2}{p{0.3\linewidth}}{
        \cellcolor[rgb]{0.9,0.9,0.9}{
            \makecell[{{p{\linewidth}}}]{
                \texttt{\tiny{[GM$|$GM]}}
                {'word': 'takso', 'assignment': 'team'}
            }
        }
    }
    & & \\ \\

    \theutterance \stepcounter{utterance}  
    & & & \multicolumn{2}{p{0.3\linewidth}}{
        \cellcolor[rgb]{0.9,0.9,0.9}{
            \makecell[{{p{\linewidth}}}]{
                \texttt{\tiny{[GM$|$GM]}}
                {'word': 'kajakas', 'assignment': 'opponent'}
            }
        }
    }
    & & \\ \\

    \theutterance \stepcounter{utterance}  
    & & & \multicolumn{2}{p{0.3\linewidth}}{
        \cellcolor[rgb]{0.9,0.9,0.9}{
            \makecell[{{p{\linewidth}}}]{
                \texttt{\tiny{[GM$|$GM]}}
                {'hidden': {'team': ['orhidee', 'jaanalind', 'jacks'], 'innocent': ['maastur', 'scotch', 'salat', 'veoauto', 'põder', 'pilkenär', 'reis'], 'opponent': ['võõrasema', 'maitseained'], 'assassin': ['varblane']}, 'revealed': {'team': {'team': ['nartsiss', 'jänes', 'rumm', 'ranne', 'peet', 'takso'], 'innocent': [], 'opponent': [], 'assassin': []}, 'opponent': {'team': [], 'innocent': [], 'opponent': ['tulp', 'traktor', 'ale', 'jalgratas', 'sibul', 'kajakas'], 'assassin': []}}}
            }
        }
    }
    & & \\ \\

    \theutterance \stepcounter{utterance}  
    & & \multicolumn{4}{p{0.6\linewidth}}{
        \cellcolor[rgb]{0.9,0.9,0.9}{
            \makecell[{{p{\linewidth}}}]{
                \texttt{\tiny{[P1$\langle$GM]}}
                \texttt{Sõna takso arvati õigesti. Teie meeskonnakaaslase käik lõppes.} \\
\texttt{Nüüd andke veel üks vihje, mis on seotud mõne teie ülejäänud meeskonnasõnaga, ja nimekiri seotud sihtmärgiks olevatest sõnadest. Ärge unustage alustada oma vihjet sõnaga 'VIHJE: ', pange uus rida ja alustage oma komadega eraldatud sihtmärgisõnade loetelu sõnaga 'SIHID: '. Märkus: mõned sõnad on võrreldes eelmiste taotlustega nimekirjadest eemaldatud.} \\
\\ 
\texttt{Teie ülejäänud meeskonnasõnad on: orhidee, jaanalind, jacks.} \\
\texttt{Teie vastase ülejäänud sõnad on: kajakas, võõrasema, maitseained.} \\
\texttt{Ülejäänud segavad sõnad on: maastur, scotch, salat, veoauto, põder, pilkenär, reis.} \\
\texttt{Ülejäänud tapmissõnad on: varblane.} \\
            }
        }
    }
    & \\ \\

    \theutterance \stepcounter{utterance}  
    & \multicolumn{4}{p{0.6\linewidth}}{
        \cellcolor[rgb]{0.9,0.9,0.9}{
            \makecell[{{p{\linewidth}}}]{
                \texttt{\tiny{[P1$\rangle$GM]}}
                \texttt{VIHJE: savann} \\
\texttt{SIHID: jaanalind} \\
            }
        }
    }
    & & \\ \\

    \theutterance \stepcounter{utterance}  
    & & & \multicolumn{2}{p{0.3\linewidth}}{
        \cellcolor[rgb]{0.9,0.9,0.9}{
            \makecell[{{p{\linewidth}}}]{
                \texttt{\tiny{[GM$|$GM]}}
                \texttt{savann} \\
            }
        }
    }
    & & \\ \\

    \theutterance \stepcounter{utterance}  
    & & & \multicolumn{2}{p{0.3\linewidth}}{
        \cellcolor[rgb]{0.9,0.9,0.9}{
            \makecell[{{p{\linewidth}}}]{
                \texttt{\tiny{[GM$|$GM]}}
                ['jaanalind']
            }
        }
    }
    & & \\ \\

    \theutterance \stepcounter{utterance}  
    & & & \multicolumn{2}{p{0.3\linewidth}}{
        \cellcolor[rgb]{0.9,0.9,0.9}{
            \makecell[{{p{\linewidth}}}]{
                \texttt{\tiny{[GM$|$GM]}}
                {'word': 'jaanalind', 'assignment': 'team'}
            }
        }
    }
    & & \\ \\

    \theutterance \stepcounter{utterance}  
    & & \multicolumn{4}{p{0.6\linewidth}}{
        \cellcolor[rgb]{0.9,0.9,0.9}{
            \makecell[{{p{\linewidth}}}]{
                \texttt{\tiny{[GM$\rangle$P2]}}
                \texttt{Sõna takso arvati õigesti. Teie käik lõppes.} \\
\texttt{Nüüd esitage veel üks komadega eraldatud loetelu vähemalt 1 ja kuni 1 sõnast järgmisest sõnade loetelust, mis on kõige paremini seotud või kõige tihedamalt seotud sõnaga 'savann'. Ärge unustage alustada oma vastust sõnadega ARVING: . Märkus: mõned sõnad on võrreldes eelmiste päringutega nimekirjast eemaldatud.} \\
\\ 
\texttt{põder, scotch, jaanalind, salat, maitseained, pilkenär, maastur, veoauto, orhidee, varblane, jacks, võõrasema, reis} \\
            }
        }
    }
    & \\ \\

    \theutterance \stepcounter{utterance}  
    & & & \multicolumn{4}{p{0.6\linewidth}}{
        \cellcolor[rgb]{0.9,0.9,0.9}{
            \makecell[{{p{\linewidth}}}]{
                \texttt{\tiny{[GM$\langle$P2]}}
                \texttt{ARVING: jaanalind} \\
            }
        }
    }
     \\ \\

    \theutterance \stepcounter{utterance}  
    & & & \multicolumn{2}{p{0.3\linewidth}}{
        \cellcolor[rgb]{0.9,0.9,0.9}{
            \makecell[{{p{\linewidth}}}]{
                \texttt{\tiny{[GM$|$GM]}}
                ['jaanalind']
            }
        }
    }
    & & \\ \\

    \theutterance \stepcounter{utterance}  
    & & & \multicolumn{2}{p{0.3\linewidth}}{
        \cellcolor[rgb]{0.9,0.9,0.9}{
            \makecell[{{p{\linewidth}}}]{
                \texttt{\tiny{[GM$|$GM]}}
                {'word': 'jaanalind', 'assignment': 'team'}
            }
        }
    }
    & & \\ \\

    \theutterance \stepcounter{utterance}  
    & & & \multicolumn{2}{p{0.3\linewidth}}{
        \cellcolor[rgb]{0.9,0.9,0.9}{
            \makecell[{{p{\linewidth}}}]{
                \texttt{\tiny{[GM$|$GM]}}
                {'word': 'jaanalind', 'assignment': 'team'}
            }
        }
    }
    & & \\ \\

    \theutterance \stepcounter{utterance}  
    & & & \multicolumn{2}{p{0.3\linewidth}}{
        \cellcolor[rgb]{0.9,0.9,0.9}{
            \makecell[{{p{\linewidth}}}]{
                \texttt{\tiny{[GM$|$GM]}}
                {'word': 'võõrasema', 'assignment': 'opponent'}
            }
        }
    }
    & & \\ \\

    \theutterance \stepcounter{utterance}  
    & & & \multicolumn{2}{p{0.3\linewidth}}{
        \cellcolor[rgb]{0.9,0.9,0.9}{
            \makecell[{{p{\linewidth}}}]{
                \texttt{\tiny{[GM$|$GM]}}
                {'hidden': {'team': ['orhidee', 'jacks'], 'innocent': ['maastur', 'scotch', 'salat', 'veoauto', 'põder', 'pilkenär', 'reis'], 'opponent': ['maitseained'], 'assassin': ['varblane']}, 'revealed': {'team': {'team': ['nartsiss', 'jänes', 'rumm', 'ranne', 'peet', 'takso', 'jaanalind'], 'innocent': [], 'opponent': [], 'assassin': []}, 'opponent': {'team': [], 'innocent': [], 'opponent': ['tulp', 'traktor', 'ale', 'jalgratas', 'sibul', 'kajakas', 'võõrasema'], 'assassin': []}}}
            }
        }
    }
    & & \\ \\

    \theutterance \stepcounter{utterance}  
    & & \multicolumn{4}{p{0.6\linewidth}}{
        \cellcolor[rgb]{0.9,0.9,0.9}{
            \makecell[{{p{\linewidth}}}]{
                \texttt{\tiny{[P1$\langle$GM]}}
                \texttt{Sõna jaanalind arvati õigesti. Teie meeskonnakaaslase käik lõppes.} \\
\texttt{Nüüd andke veel üks vihje, mis on seotud mõne teie ülejäänud meeskonnasõnaga, ja nimekiri seotud sihtmärgiks olevatest sõnadest. Ärge unustage alustada oma vihjet sõnaga 'VIHJE: ', pange uus rida ja alustage oma komadega eraldatud sihtmärgisõnade loetelu sõnaga 'SIHID: '. Märkus: mõned sõnad on võrreldes eelmiste taotlustega nimekirjadest eemaldatud.} \\
\\ 
\texttt{Teie ülejäänud meeskonnasõnad on: orhidee, jacks.} \\
\texttt{Teie vastase ülejäänud sõnad on: võõrasema, maitseained.} \\
\texttt{Ülejäänud segavad sõnad on: maastur, scotch, salat, veoauto, põder, pilkenär, reis.} \\
\texttt{Ülejäänud tapmissõnad on: varblane.} \\
            }
        }
    }
    & \\ \\

    \theutterance \stepcounter{utterance}  
    & \multicolumn{4}{p{0.6\linewidth}}{
        \cellcolor[rgb]{0.9,0.9,0.9}{
            \makecell[{{p{\linewidth}}}]{
                \texttt{\tiny{[P1$\rangle$GM]}}
                \texttt{VIHJE: epifüüt} \\
\texttt{SIHID: orhidee} \\
            }
        }
    }
    & & \\ \\

    \theutterance \stepcounter{utterance}  
    & & & \multicolumn{2}{p{0.3\linewidth}}{
        \cellcolor[rgb]{0.9,0.9,0.9}{
            \makecell[{{p{\linewidth}}}]{
                \texttt{\tiny{[GM$|$GM]}}
                \texttt{epifüüt} \\
            }
        }
    }
    & & \\ \\

    \theutterance \stepcounter{utterance}  
    & & & \multicolumn{2}{p{0.3\linewidth}}{
        \cellcolor[rgb]{0.9,0.9,0.9}{
            \makecell[{{p{\linewidth}}}]{
                \texttt{\tiny{[GM$|$GM]}}
                ['orhidee']
            }
        }
    }
    & & \\ \\

    \theutterance \stepcounter{utterance}  
    & & & \multicolumn{2}{p{0.3\linewidth}}{
        \cellcolor[rgb]{0.9,0.9,0.9}{
            \makecell[{{p{\linewidth}}}]{
                \texttt{\tiny{[GM$|$GM]}}
                {'word': 'orhidee', 'assignment': 'team'}
            }
        }
    }
    & & \\ \\

    \theutterance \stepcounter{utterance}  
    & & \multicolumn{4}{p{0.6\linewidth}}{
        \cellcolor[rgb]{0.9,0.9,0.9}{
            \makecell[{{p{\linewidth}}}]{
                \texttt{\tiny{[GM$\rangle$P2]}}
                \texttt{Sõna jaanalind arvati õigesti. Teie käik lõppes.} \\
\texttt{Nüüd esitage veel üks komadega eraldatud loetelu vähemalt 1 ja kuni 1 sõnast järgmisest sõnade loetelust, mis on kõige paremini seotud või kõige tihedamalt seotud sõnaga 'epifüüt'. Ärge unustage alustada oma vastust sõnadega ARVING: . Märkus: mõned sõnad on võrreldes eelmiste päringutega nimekirjast eemaldatud.} \\
\\ 
\texttt{põder, scotch, salat, maitseained, pilkenär, maastur, veoauto, orhidee, varblane, jacks, reis} \\
            }
        }
    }
    & \\ \\

    \theutterance \stepcounter{utterance}  
    & & & \multicolumn{4}{p{0.6\linewidth}}{
        \cellcolor[rgb]{0.9,0.9,0.9}{
            \makecell[{{p{\linewidth}}}]{
                \texttt{\tiny{[GM$\langle$P2]}}
                \texttt{ARVING: orhidee} \\
            }
        }
    }
     \\ \\

    \theutterance \stepcounter{utterance}  
    & & & \multicolumn{2}{p{0.3\linewidth}}{
        \cellcolor[rgb]{0.9,0.9,0.9}{
            \makecell[{{p{\linewidth}}}]{
                \texttt{\tiny{[GM$|$GM]}}
                ['orhidee']
            }
        }
    }
    & & \\ \\

    \theutterance \stepcounter{utterance}  
    & & & \multicolumn{2}{p{0.3\linewidth}}{
        \cellcolor[rgb]{0.9,0.9,0.9}{
            \makecell[{{p{\linewidth}}}]{
                \texttt{\tiny{[GM$|$GM]}}
                {'word': 'orhidee', 'assignment': 'team'}
            }
        }
    }
    & & \\ \\

    \theutterance \stepcounter{utterance}  
    & & & \multicolumn{2}{p{0.3\linewidth}}{
        \cellcolor[rgb]{0.9,0.9,0.9}{
            \makecell[{{p{\linewidth}}}]{
                \texttt{\tiny{[GM$|$GM]}}
                {'word': 'orhidee', 'assignment': 'team'}
            }
        }
    }
    & & \\ \\

    \theutterance \stepcounter{utterance}  
    & & & \multicolumn{2}{p{0.3\linewidth}}{
        \cellcolor[rgb]{0.9,0.9,0.9}{
            \makecell[{{p{\linewidth}}}]{
                \texttt{\tiny{[GM$|$GM]}}
                {'word': 'maitseained', 'assignment': 'opponent'}
            }
        }
    }
    & & \\ \\

    \theutterance \stepcounter{utterance}  
    & & & \multicolumn{2}{p{0.3\linewidth}}{
        \cellcolor[rgb]{0.9,0.9,0.9}{
            \makecell[{{p{\linewidth}}}]{
                \texttt{\tiny{[GM$|$GM]}}
                {'hidden': {'team': ['jacks'], 'innocent': ['maastur', 'scotch', 'salat', 'veoauto', 'põder', 'pilkenär', 'reis'], 'opponent': [], 'assassin': ['varblane']}, 'revealed': {'team': {'team': ['nartsiss', 'jänes', 'rumm', 'ranne', 'peet', 'takso', 'jaanalind', 'orhidee'], 'innocent': [], 'opponent': [], 'assassin': []}, 'opponent': {'team': [], 'innocent': [], 'opponent': ['tulp', 'traktor', 'ale', 'jalgratas', 'sibul', 'kajakas', 'võõrasema', 'maitseained'], 'assassin': []}}}
            }
        }
    }
    & & \\ \\

    \theutterance \stepcounter{utterance}  
    & & \multicolumn{4}{p{0.6\linewidth}}{
        \cellcolor[rgb]{0.9,0.9,0.9}{
            \makecell[{{p{\linewidth}}}]{
                \texttt{\tiny{[P1$\langle$GM]}}
                \texttt{Sõna orhidee arvati õigesti. Teie meeskonnakaaslase käik lõppes.} \\
\texttt{Nüüd andke veel üks vihje, mis on seotud mõne teie ülejäänud meeskonnasõnaga, ja nimekiri seotud sihtmärgiks olevatest sõnadest. Ärge unustage alustada oma vihjet sõnaga 'VIHJE: ', pange uus rida ja alustage oma komadega eraldatud sihtmärgisõnade loetelu sõnaga 'SIHID: '. Märkus: mõned sõnad on võrreldes eelmiste taotlustega nimekirjadest eemaldatud.} \\
\\ 
\texttt{Teie ülejäänud meeskonnasõnad on: jacks.} \\
\texttt{Teie vastase ülejäänud sõnad on: maitseained.} \\
\texttt{Ülejäänud segavad sõnad on: maastur, scotch, salat, veoauto, põder, pilkenär, reis.} \\
\texttt{Ülejäänud tapmissõnad on: varblane.} \\
            }
        }
    }
    & \\ \\

    \theutterance \stepcounter{utterance}  
    & \multicolumn{4}{p{0.6\linewidth}}{
        \cellcolor[rgb]{0.9,0.9,0.9}{
            \makecell[{{p{\linewidth}}}]{
                \texttt{\tiny{[P1$\rangle$GM]}}
                \texttt{VIHJE: valetid} \\
\texttt{SIHID: jacks} \\
            }
        }
    }
    & & \\ \\

    \theutterance \stepcounter{utterance}  
    & & & \multicolumn{2}{p{0.3\linewidth}}{
        \cellcolor[rgb]{0.9,0.9,0.9}{
            \makecell[{{p{\linewidth}}}]{
                \texttt{\tiny{[GM$|$GM]}}
                \texttt{valetid} \\
            }
        }
    }
    & & \\ \\

    \theutterance \stepcounter{utterance}  
    & & & \multicolumn{2}{p{0.3\linewidth}}{
        \cellcolor[rgb]{0.9,0.9,0.9}{
            \makecell[{{p{\linewidth}}}]{
                \texttt{\tiny{[GM$|$GM]}}
                ['jacks']
            }
        }
    }
    & & \\ \\

    \theutterance \stepcounter{utterance}  
    & & & \multicolumn{2}{p{0.3\linewidth}}{
        \cellcolor[rgb]{0.9,0.9,0.9}{
            \makecell[{{p{\linewidth}}}]{
                \texttt{\tiny{[GM$|$GM]}}
                {'word': 'jacks', 'assignment': 'team'}
            }
        }
    }
    & & \\ \\

    \theutterance \stepcounter{utterance}  
    & & & \multicolumn{2}{p{0.3\linewidth}}{
        \cellcolor[rgb]{0.9,0.9,0.9}{
            \makecell[{{p{\linewidth}}}]{
                \texttt{\tiny{[GM$|$GM]}}
                \texttt{opponent has won} \\
            }
        }
    }
    & & \\ \\

\end{supertabular}
}

\end{document}
